\section{Порождающие и знак перестановки}
\begin{definition}
\term{Подгруппа} (\emph{subgroup}) группы \(G\) --- подмножество \(H\), которое содержит единицу и замкнуто относительно умножения и взятия обратного. С тем же действием оно само является группой.
\end{definition}
\begin{definition}
\term{Подгруппа, порождённая множеством} \(S\subseteq G\) (\emph{generated subgroup}), обозначается \(\langle S\rangle\) и состоит из всех конечных произведений элементов \(S\) и их обратных; пустое произведение равно единице. \par\term{Конечно порождённая группа} (\emph{finitely generated group}) --- группа, для которой \(G=\langle S\rangle\) для некоторого конечного \(S\).
\end{definition}
\begin{definition}
\term{Инверсия перестановки} (\emph{inversion}) \(\sigma\in S_n\) --- пара индексов \(i<j\), для которой \(\sigma(i)>\sigma(j)\). \par\term{Знак перестановки} (\emph{sign of a permutation}), имеющей \(N(\sigma)\) инверсий, равен \(\operatorname{sgn}\sigma=(-1)^{N(\sigma)}\). Это чётность перестановки, которую не следует смешивать с чётностью её порядка как элемента группы.
\end{definition}
\begin{theorem}[Мультипликативность знака]\label{thm:olympiad-permutation-sign}
Для перестановок \(\sigma,\tau\in S_n\) выполнено \(\operatorname{sgn}(\sigma\circ\tau)=(\operatorname{sgn}\sigma)(\operatorname{sgn}\tau)\). Знак любой транспозиции равен \(-1\).
\end{theorem}
\begin{proof}
Обмен соседних позиций меняет число инверсий ровно на один: меняется порядок этой пары, а суммарные отношения с прочими позициями сохраняются. Любую перестановку можно отсортировать соседними обменами, каждый раз убирая одну инверсию. Поэтому чётность числа соседних обменов в любом её разложении равна \(N(\sigma)\) по модулю \(2\). При композиции разложения приписываются друг к другу, и их длины складываются; отсюда первая формула. Обмен позиций \(i<j\) можно выполнить \(2(j-i)-1\) соседними обменами; это нечётное число, что доказывает вторую формулу.
\end{proof}
